% SPDX-License-Identifier: CC-BY-4.0
% Copyright (C) 2026 the Langlandscape team
% Status-tag legend (round R5). Shared: start of chapter s00 (main.tex), after the abstract of the article and of the
% submission report. No \input; uses only the status macros of preamble.tex / submission/src/macros.tex.
\begin{sloppypar}
\noindent\textbf{Status tags.} Every result carries its status as a tag in brackets, printed exactly as follows.
\begin{itemize}\setlength{\itemsep}{0pt}
\item \stproved: the proof is written out in our chapters.
\item \stpmc: proved, using the named cited results as black boxes; \status{proved modulo cited geometry}: the same, where
  the cited inputs are the geometric results named there (Deligne--Rapoport, Deuring, Grothendieck's character group).
\item \stcited: taken from the literature with a precise pointer, not reproved; \stknown: standard or textbook, with a
  reference.
\item \status{definition}: a definition, nothing to prove; \sthyp: a named hypothesis, assumed only where stated and never
  claimed.
\item \stcomp: checked by machine computation over the stated range; a computation is never part of a proof.
\item \stformal: a Lean~4 declaration with this content compiles with the standard axioms only, named next to it as
  \textup{\small(Lean: \dots)}; \emph{formalized modulo $H$} (\emph{Lean-conditional}): the same, with the registered
  hypotheses $H$ as explicit arguments; \textup{\small(Lean: planned)}: not yet formalized. In the claims table of the
  submission report, \emph{Lean-verified} means formalized without qualifier, with no difference recorded.
\item \stsketch, \stexp, \stconj, \stopen, \sttoprove: no proof is claimed (a sketch is given; expected to hold; a
  conjecture; an open question; a proof still to be written).
\item Notes on changes after internal audits are printed as \textup{\textsf{\small[provisional: \dots]}} in the chapters and
  as \textup{\textsf{\small[revision note: \dots]}} in the submission report.
\end{itemize}
\end{sloppypar}
